% Inactive research fragment. These examples check the scope of a possible
% comparison argument; they do not disprove a domination theorem.
% Setting: a neighborhood of the closed unit ball in C^3, m=2.

Let $\beta_j=\ddc|z_j|^2$ and $\beta=\beta_1+\beta_2+\beta_3$.
We distinguish the shifted Hessian measure
$H_2^\beta(t)=(\beta+\ddc t)^2\wedge\beta$ from the raw expression
$L_2^\beta(t)=(\ddc t)^2\wedge\beta$.

\begin{lemma}\label{draft635:shifted-raw-order}
There are smooth bounded $\beta$-plurisubharmonic functions $u,v$ on
the unit ball such that
\[
 H_2^\beta(v)>H_2^\beta(u),\qquad
 L_2^\beta(v)<L_2^\beta(u).
\]
Neither function is subharmonic for the unshifted $2$-Hessian cone
relative to $\beta$.
\end{lemma}

\begin{proof}
Take $u=-\tfrac12\sum_{j=1}^3|z_j|^2$ and
$v=-\tfrac45|z_1|^2$.
The eigenvalues of $\beta+\ddc u$ relative to $\beta$ are
$(\tfrac12,\tfrac12,\tfrac12)$, whereas those of
$\beta+\ddc v$ are $(\tfrac15,1,1)$.
Both forms are strictly positive. For a diagonal real $(1,1)$-form
$A$ with eigenvalues $\lambda_1,\lambda_2,\lambda_3$,
\[
 A^2\wedge\beta
 =\frac{\lambda_1\lambda_2+\lambda_1\lambda_3+
         \lambda_2\lambda_3}{3}\,\beta^3.
\]
Consequently,
\[
 H_2^\beta(u)=\tfrac14\beta^3,
 \qquad H_2^\beta(v)=\tfrac7{15}\beta^3,
 \qquad L_2^\beta(u)=\tfrac14\beta^3,
 \qquad L_2^\beta(v)=0.
\]
The traces of $\ddc u$ and $\ddc v$ relative to $\beta$ are
$-\tfrac32$ and $-\tfrac45$, respectively. They therefore lie
outside the unshifted $2$-Hessian cone.
\end{proof}

This example concerns an implication between two operator inequalities.
It has no prescribed boundary order and is not a counterexample to
a boundary comparison principle. In this closed constant background,
adding $\sum_j|z_j|^2$ to both functions gives a common local potential
correction. Such an exact correction requires $\ddc\rho=\beta$;
a nonclosed Hermitian background cannot have an exact correction
$\ddc\rho=\omega$ on a neighborhood where $d\omega\ne0$.
The book's correction with $\ddc\rho\geq2\omega$ instead leaves the
smooth remainder $\omega-\ddc\rho$ in the shifted products.

\begin{lemma}\label{draft635:admissible-not-positive}
There is a smooth bounded shifted $2$-admissible function $v\geq1$
on the closed unit ball such that the following construction remains
$2$-admissible but its perturbed form is not semipositive:
\[
 a=\tfrac9{10},\qquad
 \rho=\tfrac1{100}\left(\sum_{j=1}^3|z_j|^2-1\right),
 \qquad \varphi=av-(1-a)\rho.
\]
Here $\rho=0$ on the boundary, $\rho<0$ in the interior,
$|\rho|\leq v$ and $0<\ddc\rho\leq\beta$.
\end{lemma}

\begin{proof}
Set $v=3-\tfrac75|z_1|^2$. The relative eigenvalues of
$\beta+\ddc v$ are $(-\tfrac25,1,1)$.
Their first and second symmetric polynomials are $\tfrac85$ and
$\tfrac15$, both strictly positive, and $v\geq\tfrac85$ on the
closed unit ball. Thus $v$ is shifted $2$-admissible but is not
$\beta$-plurisubharmonic.
The relative eigenvalues of $\beta+\ddc\varphi$ are
\[
 \left(-\tfrac{261}{1000},\tfrac{999}{1000},
       \tfrac{999}{1000}\right).
\]
Their first and second symmetric polynomials are
$\tfrac{1737}{1000}$ and $\tfrac{476523}{1000000}$.
The form therefore belongs to the open $2$-Hessian cone and has
a negative eigenvalue.
\end{proof}

Cone convexity still gives
\[
 \beta+\ddc\varphi
 =a(\beta+\ddc v)+(1-a)(\beta-\ddc\rho)
 \in\Gamma_2(\beta).
\]
Thus an argument stated for plurisubharmonic functions can retain
this construction in the larger admissible cone, but its assertion
of semipositivity must be replaced by the cone statement. A separate
comparison theorem for the shifted operator, or a proof accounting
for the smooth remainder after correction, is still needed before
using an unshifted local comparison theorem.
